\subsection{3.4. Tìm kiếm A*}

\subsubsection{3.4.1. Ý tưởng}

A* {[}7, §3.5.2{]} mở rộng nút có \(f(n) = g(n) + h(n)\) nhỏ nhất, trong đó \(g(n)\) là tổng chi phí từ trạng thái đầu tới \(n\), còn \(h(n)\) là ước lượng chi phí còn lại tới đích. Khi \(h \equiv 0\), A* chính là UCS. Khi \(h\) chấp nhận được (không bao giờ ước lượng quá chi phí thật), A* trả về lời giải tối ưu; khi \(h\) còn nhất quán, mỗi trạng thái chỉ cần mở rộng một lần. Heuristic càng sát chi phí thật thì A* càng mở rộng ít nút hơn UCS.

Trạng thái, hành động, hàm chuyển và kiểm tra đích dùng chung với ba thuật toán còn lại (chương 2). A* dùng cùng hàm chi phí với UCS (mục 3.3), để so được tổng chi phí của hai thuật toán.

\subsubsection{3.4.2. Mã giả}

\begin{verbatim}
A*(problem, h):
    s0 ← problem.initial_state()
    nếu h(s0) = ∞: trả về NO_SOLUTION                  ▷ ngõ cụt ngay từ đầu
    frontier ← hàng đợi ưu tiên chứa (f = h(s0), h(s0), thứ_tự, g = 0, s0)
    best_g[s0] ← 0
    lặp khi frontier khác rỗng:
        (f, h_s, _, g, s) ← lấy ra mục nhỏ nhất theo (f, h, thứ_tự)
        nếu g > best_g[s]: bỏ qua                        ▷ mục cũ (lazy deletion)
        nếu problem.is_goal(s): trả về đường đi tới s, chi phí g
        expanded ← expanded + 1
        với mỗi a trong problem.actions(s):               ▷ thứ tự U, D, L, R (đã tách: U1 … R2)
            s' ← problem.result(s, a)
            nếu s' = None: bỏ qua                         ▷ khối rơi, không phải nút con
            generated ← generated + 1
            g' ← g + problem.step_cost(s, a, s')
            nếu g' ≥ best_g[s'] (nếu có): bỏ qua
            best_g[s'] ← g';  cha[s'] ← (s, a)
            nếu h(s') = ∞: bỏ qua                         ▷ ngõ cụt, không đưa vào biên
            đưa (g' + h(s'), h(s'), thứ_tự mới, g', s') vào frontier
    trả về NO_SOLUTION
\end{verbatim}

Mã nguồn: \texttt{Source/\allowbreak{}solvers/\allowbreak{}astar.\allowbreak{}py} (lớp \texttt{AStar}, \texttt{AStarM}) và \texttt{Source/\allowbreak{}solvers/\allowbreak{}heuristics.\allowbreak{}py}.

\subsubsection{3.4.3. Hai heuristic}

Gọi \(D(x)\) là một khoảng cách từ ô \(x\) tới đích (định nghĩa riêng cho từng heuristic bên dưới), và với mỗi trạng thái \(s\), gọi \(\Phi(s) = \min_{x \in \text{cells}(s)} D(x)\), trong đó \(\text{cells}(s)\) là các ô khối đang chiếm: một ô khi khối đứng, hai ô khi khối nằm, hai ô của hai khối con khi đã tách. Cả hai heuristic có chung dạng

\[
h(s) = \begin{cases}
0 & \text{nếu } s \text{ là trạng thái đích,} \\
\infty & \text{nếu có } x \in \text{cells}(s) \text{ với } D(x) = \infty, \\
\max\left(\left\lceil \Phi(s)/2 \right\rceil,\ 1\right) & \text{còn lại.}
\end{cases}
\]

Ý nghĩa: một lần lăn, khối nguyên tiến lại gần đích nhiều nhất 2 ô, khối con nhiều nhất 1 ô, nên cần ít nhất \(\lceil \Phi/2 \rceil\) lần lăn nữa. Trạng thái chưa phải đích thì còn ít nhất một bước, nên lấy max với 1.

\textbf{Bẫy công tắc tách.} Công tắc tách đưa hai khối con tới hai ô đáp ở xa. Nếu \(D\) là khoảng cách Manhattan thuần tới đích, heuristic có thể ước lượng quá. Ví dụ (dùng làm test trong \texttt{tests/\allowbreak{}test\_\allowbreak{}astar.\allowbreak{}py}): khối đứng ở (5, 0), công tắc tách ở (5, 3), hai ô đáp (1, 10) và (1, 12), đích (1, 13). Lời giải thật chỉ 4 bước (R, R, R1 rồi tự ghép, R), nhưng Manhattan thuần cho \(\lceil 17/2 \rceil = 9 > 4\). Nhóm khóa trước {[}8{]} mắc đúng lỗi này. Cả hai heuristic của nhóm đều tính đường tắt qua công tắc tách:

\begin{itemize}
\tightlist
\item
  \textbf{\(h_M\) (Manhattan, khóa \texttt{astar\_\allowbreak{}m}).} \(D_M(x) = \min_{t \in T} |x_r - t_r| + |x_c - t_c|\), với \(T\) gồm ô đích và mọi công tắc tách. Không cần tiền xử lý. Ở màn không có công tắc tách, đây là Manhattan thường. Trong ví dụ trên, \(h_M = \lceil 3/2 \rceil = 2\).
\item
  \textbf{\(h_P\) (lưới điểm, khóa \texttt{astar}, heuristic chính).} Dựng đồ thị trên các ô không trống (kể cả cầu đang đóng, ô dễ vỡ, mọi loại công tắc): hai ô kề cạnh nối bằng cạnh chi phí 1 theo cả hai chiều; mỗi công tắc tách \(X\) có thêm cạnh chi phí 0 từ \(X\) tới hai ô đáp \(A\), \(B\) của nó. \(D_P(x)\) là độ dài đường đi ngắn nhất từ \(x\) tới đích trên đồ thị này, tính trước cho mọi ô trong \texttt{prepare()} bằng 0-1 BFS ngược từ đích, độ phức tạp \(O(R \cdot C)\) với bàn \(R \times C\). Ô không có đường tới đích có \(D_P = \infty\). Trong ví dụ trên, \(D_P(5, 3) = \min(D_P(1, 10), D_P(1, 12)) = 1\), nên \(h_P = \lceil 4/2 \rceil = 2\).
\end{itemize}

\(h_P\) tính cả tường và hố (ô trống), nên luôn có \(D_P \ge D_M\) trên mọi ô tới được đích: \(h_P\) mạnh hơn hoặc bằng \(h_M\) (Bảng 3.4.2). \(h_P\) còn phát hiện ngõ cụt: khối hoặc một khối con nằm ở vùng sàn không nối được với đích thì \(h_P = \infty\) và A* bỏ trạng thái đó ngay.

\subsubsection{3.4.4. Chứng minh chấp nhận được và nhất quán}

\textbf{Giả thiết.} Mọi chi phí bước \(c(s, a, s') \ge 1\). Hàm chi phí của nhóm thỏa giả thiết này: mọi bước tốn ít nhất 1,10 (mục 3.3.2). Chứng minh dưới đây chỉ dùng giả thiết đó, nên vẫn đúng nếu đổi các mức chi phí, miễn là mỗi bước tốn ít nhất 1.

\textbf{Hai tính chất của \(D\)} (đúng cho cả \(D_M\) và \(D_P\)):

\begin{itemize}
\tightlist
\item
  \begin{enumerate}
  \def\labelenumi{(\Alph{enumi})}
  \setcounter{enumi}{11}
  \tightlist
  \item
    Hai ô không trống kề cạnh \(u, v\) thì \(D(u) \le D(v) + 1\). Với \(D_M\): Manhattan tới một tập điểm thay đổi tối đa 1 khi dịch 1 ô. Với \(D_P\): \(u, v\) nối bằng cạnh chi phí 1, nên đi từ \(u\) sang \(v\) rồi theo đường ngắn nhất của \(v\).
  \end{enumerate}
\item
  \begin{enumerate}
  \def\labelenumi{(\Alph{enumi})}
  \setcounter{enumi}{18}
  \tightlist
  \item
    Công tắc tách \(X\) với ô đáp \(A, B\) thì \(D(X) \le \min(D(A), D(B))\). Với \(D_M\): \(X \in T\) nên \(D_M(X) = 0\). Với \(D_P\): có cạnh chi phí 0 từ \(X\) tới \(A\) và tới \(B\).
  \end{enumerate}
\end{itemize}

\textbf{Bổ đề 1.} Với mọi bước chuyển hợp lệ \(s \to s'\): \(\Phi(s) \le \Phi(s') + 2\) nếu khối nguyên lăn, và \(\Phi(s) \le \Phi(s') + 1\) nếu một khối con đi.

\emph{Chứng minh.} Xét bước lăn của khối nguyên, chưa tính dịch chuyển. Theo bảng chuyển (Bảng 2.4), mỗi ô mới của khối hoặc kề cạnh một ô cũ, hoặc kề cạnh một ô mới khác mà ô đó kề cạnh một ô cũ; mọi ô này đều không trống vì bước chuyển hợp lệ. Gọi \(x'\) là ô mới có \(D(x') = \Phi(s')\); theo (L), có ô cũ \(x\) với \(D(x) \le D(x') + 2\), nên \(\Phi(s) \le D(x) \le \Phi(s') + 2\). Khối con đi 1 ô thì ô mới kề ô cũ, khối con kia đứng yên: tương tự với \(+1\). Ghép hai khối con không đổi ô nào, nên \(\Phi\) không đổi. Đổi cầu cũng không đổi ô nào.

Còn lại bước có dịch chuyển: khối lăn tới trạng thái trung gian \(m\) đứng trên công tắc tách \(X\), rồi thành \(s'\) với hai khối con ở \(A, B\). Theo phần trên, \(\Phi(s) \le \Phi(m) + 2 = D(X) + 2\), và theo (S), \(D(X) \le \min(D(A), D(B)) = \Phi(s')\). Vậy \(\Phi(s) \le \Phi(s') + 2\). ∎

\textbf{Định lý 1 (nhất quán).} Với mọi bước chuyển hợp lệ \(s \to s'\) mà \(h(s) < \infty\): \(h(s) \le c(s, a, s') + h(s')\).

\emph{Chứng minh.} Nếu \(h(s') = \infty\) thì hiển nhiên. Nếu \(s\) là đích thì \(h(s) = 0\). Còn lại, \(h(s) = \max(\lceil \Phi(s)/2 \rceil, 1)\).

\begin{itemize}
\tightlist
\item
  \(s'\) là đích: \(\Phi(s') = 0\) (ô đích có \(D = 0\)), theo Bổ đề 1, \(\Phi(s) \le 2\), nên \(h(s) = 1 \le c(s, a, s') = c + h(s')\).
\item
  \(s'\) không phải đích: theo Bổ đề 1, \(\lceil \Phi(s)/2 \rceil \le \lceil (\Phi(s') + 2)/2 \rceil = \lceil \Phi(s')/2 \rceil + 1 \le h(s') + 1\); và \(1 \le h(s') + 1\). Vì \(c \ge 1\), ta có \(h(s) \le h(s') + 1 \le h(s') + c\). ∎
\end{itemize}

\textbf{Bổ đề 2 (ngõ cụt).} Nếu \(h_P(s) = \infty\) thì từ \(s\) không tới được đích.

\emph{Chứng minh.} Mọi chuyển động thật của khối đều đi theo cạnh của đồ thị lưới điểm: lăn và đi của khối con theo các ô kề, dịch chuyển theo cạnh \(X \to A\), \(X \to B\). Vì vậy mọi ô mà một phần khối chiếm về sau đều đi tới được, trên đồ thị, từ một ô khối đang chiếm. Nếu khối nguyên có một ô với \(D_P = \infty\) thì ô kia kề nó cũng vậy (cạnh lưới đi được hai chiều). Nếu đã tách và khối con ở ô \(u\) có \(D_P(u) = \infty\), thì mọi vị trí \(u'\) sau này của khối con này cũng có \(D_P(u') = \infty\); khi hai khối con ghép tại \(u'\) và \(v'\) kề nhau, \(D_P(v') = \infty\) theo (L). Trong mọi trường hợp, khối không bao giờ chiếm ô đích, nên không thắng được. ∎

\textbf{Hệ quả (chấp nhận được và tối ưu).} Với mọi trạng thái \(s\), gọi \(h^*(s)\) là chi phí thật nhỏ nhất từ \(s\) tới đích. Nếu \(h^*(s) = \infty\) thì bất đẳng thức hiển nhiên. Nếu \(h^*(s) < \infty\), theo Bổ đề 2 thì \(h(s) < \infty\) (với \(h_M\) thì luôn hữu hạn), và cộng Định lý 1 dọc một lời giải tối ưu \(s = s_0, s_1, \dots, s_k\) (đích) với \(h(s_k) = 0\) cho \(h(s) \le h^*(s)\). Vậy \(h_M\) và \(h_P\) đều chấp nhận được và nhất quán, nên A* trong mục 3.4.2 (kiểm tra đích khi lấy ra, giữ \texttt{best\_\allowbreak{}g}) trả về lời giải có tổng chi phí nhỏ nhất {[}7, §3.5.2{]}, bằng UCS, và việc bỏ các trạng thái có \(h = \infty\) không làm mất lời giải.

\textbf{Không được phạt thêm.} Đề gợi ý cộng thêm một phạt nhỏ khi khối chưa đứng, vì cần ít nhất một lần lăn nữa để đứng lên đích {[}1, mục 3.5{]}. Với dạng \(\lceil \Phi/2 \rceil\) của nhóm, phạt này làm mất tính chấp nhận được, nên nhóm không dùng. Ví dụ \(\lceil \Phi/2 \rceil + 1\) khi khối đang nằm là \textbf{không} chấp nhận được: khối nằm sát đích chỉ cần 1 bước để đứng lên đích, trong khi heuristic cho 2.

\subsubsection{3.4.5. Kiểm tra bằng số liệu}

Bổ sung cho chứng minh, công cụ \texttt{Source/\allowbreak{}tools/\allowbreak{}check\_\allowbreak{}heuristics.\allowbreak{}py} duyệt toàn bộ không gian trạng thái tới được của từng màn và đếm:

\begin{itemize}
\tightlist
\item
  vi phạm nhất quán: cạnh \((s, a, s')\) có \(h(s) > c(s, a, s') + h(s')\);
\item
  vi phạm chấp nhận được: trạng thái trên một lời giải tối ưu có \(h(s)\) lớn hơn chi phí thật còn lại;
\item
  màn mà A* có tổng chi phí khác UCS, hoặc mở rộng nhiều nút hơn UCS.
\end{itemize}

\begin{longtable}[]{@{}lll@{}}
\caption{\textbf{Bảng 3.4.1.} Kết quả kiểm tra trên 33 màn gốc, với hàm chi phí ở Bảng 3.3.1.}\tabularnewline
\toprule\noalign{}
\textbf{Hạng mục} & \textbf{\(h_M\)} & \textbf{\(h_P\)} \\
\midrule\noalign{}
\endfirsthead
\toprule\noalign{}
\textbf{Hạng mục} & \textbf{\(h_M\)} & \textbf{\(h_P\)} \\
\midrule\noalign{}
\endhead
\bottomrule\noalign{}
\endlastfoot
Số trạng thái tới được (tổng 33 màn) & 70 062 & 70 062 \\
Số cạnh đã kiểm tra & 340 414 & 340 414 \\
Vi phạm nhất quán & 0 & 0 \\
Vi phạm chấp nhận được (dọc lời giải tối ưu) & 0 & 0 \\
Màn có tổng chi phí khác UCS & 0 / 33 & 0 / 33 \\
Màn mở rộng nhiều nút hơn UCS & 0 / 33 & 0 / 33 \\
Màn có số bước bằng số bước tối ưu tham khảo & 33 / 33 & 33 / 33 \\
\end{longtable}

\begin{longtable}[]{@{}
  >{\raggedright\arraybackslash}p{(\linewidth - 6\tabcolsep) * \real{0.4028}}
  >{\raggedright\arraybackslash}p{(\linewidth - 6\tabcolsep) * \real{0.1250}}
  >{\raggedright\arraybackslash}p{(\linewidth - 6\tabcolsep) * \real{0.2361}}
  >{\raggedright\arraybackslash}p{(\linewidth - 6\tabcolsep) * \real{0.2361}}@{}}
\caption{\textbf{Bảng 3.4.2.} Tổng số nút mở rộng và thời gian tìm kiếm trên 33 màn gốc (lần đo ở mục 5.3, hàm chi phí ở Bảng 3.3.1; thời gian là trung vị 5 lần đo).}\tabularnewline
\toprule\noalign{}
\begin{minipage}[b]{\linewidth}\raggedright
\end{minipage} & \begin{minipage}[b]{\linewidth}\raggedright
\textbf{UCS}
\end{minipage} & \begin{minipage}[b]{\linewidth}\raggedright
\textbf{A* với \(h_M\)}
\end{minipage} & \begin{minipage}[b]{\linewidth}\raggedright
\textbf{A* với \(h_P\)}
\end{minipage} \\
\midrule\noalign{}
\endfirsthead
\toprule\noalign{}
\begin{minipage}[b]{\linewidth}\raggedright
\end{minipage} & \begin{minipage}[b]{\linewidth}\raggedright
\textbf{UCS}
\end{minipage} & \begin{minipage}[b]{\linewidth}\raggedright
\textbf{A* với \(h_M\)}
\end{minipage} & \begin{minipage}[b]{\linewidth}\raggedright
\textbf{A* với \(h_P\)}
\end{minipage} \\
\midrule\noalign{}
\endhead
\bottomrule\noalign{}
\endlastfoot
Tổng số nút mở rộng & 35 281 & 33 474 (−5,1 \%) & 32 712 (−7,3 \%) \\
Trung vị phần trăm nút bớt được so với UCS, theo màn & & 4,3 \% & 5,7 \% \\
Màn bớt nhiều nhất & & màn 1: −60,6 \% & màn 1: −60,6 \% \\
Tổng thời gian tìm kiếm (ms) & 1 163,2 & 1 193,3 & 1 130,6 \\
Thời gian trung vị cho mỗi nút mở rộng (µs) & 18,6 & 19,9 & 19,2 \\
\end{longtable}

Ngoài 33 màn gốc, công cụ còn kiểm 2 015 147 cạnh của 14 màn tự sinh (mục 5.2), cũng không có vi phạm nào. Kết luận chính: cả hai heuristic đều đúng (0 vi phạm, cùng chi phí với UCS ở 33/33 màn); \(h_P\) bớt nhiều nút hơn \(h_M\) và rẻ hơn khi tính cho mỗi nút (tra bảng thay vì tính Manhattan tới nhiều đích). Nhận xét chi tiết ở mục 5.4.4.

\subsubsection{3.4.6. Độ phức tạp, tính đầy đủ và tối ưu}

\begin{itemize}
\tightlist
\item
  \textbf{Đầy đủ.} Không gian trạng thái hữu hạn (vị trí và tư thế của khối, hoặc vị trí hai khối con, nhân với cấu hình cầu), mọi chi phí \(\ge 1\), và A* giữ \texttt{best\_\allowbreak{}g} nên không lặp vô hạn. A* luôn dừng: hoặc tìm ra lời giải, hoặc trả về NO\_SOLUTION khi biên rỗng.
\item
  \textbf{Tối ưu.} Theo Hệ quả ở mục 3.4.4.
\item
  \textbf{Thời gian.} Trường hợp xấu nhất như UCS: \(O(b^{1 + \lfloor C^*/\varepsilon \rfloor})\) với \(b \le 8\) (4 hướng khi khối nguyên, 8 hành động khi đã tách), \(C^*\) là chi phí tối ưu, \(\varepsilon = 1{,}10\) là chi phí bước nhỏ nhất {[}7, §3.4.2{]}. Vì có tập đã thăm, số trạng thái mở rộng cũng bị chặn bởi \(|S|\), mỗi thao tác trên heap tốn \(O(\log |S|)\). Heuristic tốt làm giảm số nút mở rộng trên thực tế.
\item
  \textbf{Bộ nhớ.} \(O(|S|)\) cho biên, \texttt{best\_\allowbreak{}g} và bảng cha; \(h_P\) thêm bảng khoảng cách \(O(R \cdot C)\).
\end{itemize}

\subsubsection{3.4.7. Các điểm riêng khi cài đặt}

\begin{itemize}
\tightlist
\item
  \textbf{Phá hòa.} Khi \(f\) bằng nhau, ưu tiên \(h\) nhỏ hơn (gần đích hơn), rồi theo thứ tự đưa vào biên. Mục trong heap là \texttt{(f, h, thứ\_\allowbreak{}tự, g, s)}; \texttt{thứ\_\allowbreak{}tự} là bộ đếm tăng dần, để Python không phải so sánh hai trạng thái.
\item
  \textbf{Lazy deletion.} Khi tìm được \(g\) nhỏ hơn cho một trạng thái đã có trong biên, A* đưa thêm một mục mới thay vì sửa mục cũ; lúc lấy ra, mục có \(g >\) \texttt{best\_\allowbreak{}g} bị bỏ qua và không tính vào số nút mở rộng. Vì vậy \texttt{max\_\allowbreak{}frontier} tính cả các mục cũ.
\item
  \textbf{Kiểm tra đích khi lấy ra} khỏi hàng đợi, không phải khi sinh nút, để giữ tính tối ưu.
\item
  \textbf{Cắt ngõ cụt.} Trạng thái có \(h = \infty\) được ghi vào \texttt{best\_\allowbreak{}g} (tính là đã thăm) nhưng không đưa vào biên. Ở màn không có lời giải, \(h_P\) có thể kết luận ngay từ trạng thái đầu mà không mở rộng nút nào.
\item
  \textbf{Heuristic dựng trong \texttt{prepare()}} chứ không trong \texttt{search()}, để thực nghiệm tách được thời gian chuẩn bị khỏi thời gian tìm kiếm (mục 5.1).
\item
  \textbf{Hai biến thể dùng chung code.} \texttt{AStarM} là lớp con của \texttt{AStar}, chỉ khác thuộc tính \texttt{heuristic}. GUI dùng \texttt{astar} (\(h_P\)) cho nút Giải bằng A*, và chạy cả hai trong nút So sánh (mục 4.6).
\item
  \textbf{Kiểm thử.} \texttt{tests/\allowbreak{}test\_\allowbreak{}solvers.\allowbreak{}py} (mê cung mẫu, có ô chi phí khác 1), \texttt{tests/\allowbreak{}test\_\allowbreak{}astar.\allowbreak{}py} (bẫy công tắc tách, ngõ cụt, khớp số bước tối ưu, 0 vi phạm) và \texttt{tools/\allowbreak{}check\_\allowbreak{}heuristics.\allowbreak{}py} (đủ 33 màn).
\end{itemize}
